\documentclass[11pt,a4paper]{article}

% ---------------------------------------------------------------------------
% FSB Nexus - Project Blueprint
% Compile with: pdflatex FSB_Nexus_Blueprint.tex
% ---------------------------------------------------------------------------

\usepackage[utf8]{inputenc}
\usepackage[T1]{fontenc}
\usepackage[english]{babel}
\usepackage{lmodern}
\usepackage[margin=2.5cm]{geometry}
\usepackage{enumitem}
\usepackage{titlesec}
\usepackage{xcolor}
\usepackage{booktabs}
\usepackage{longtable}
\usepackage{parskip}
\usepackage[hidelinks]{hyperref}

\definecolor{fsbblue}{RGB}{20,60,110}

\titleformat{\section}{\large\bfseries\color{fsbblue}}{\thesection}{0.6em}{}
\titleformat{\subsection}{\normalsize\bfseries}{\thesubsection}{0.6em}{}

\setlist[itemize]{leftmargin=1.4em,itemsep=2pt,topsep=3pt}

\title{\vspace{-1.2cm}\textbf{\color{fsbblue}FSB Nexus}\\[0.3em]
\large An AI-Powered Digital Campus Assistant for the\\Faculty of Sciences of Bizerte\\[0.6em]
\normalsize Project Blueprint and Roadmap}

\author{}
\date{}

\begin{document}

\maketitle
\vspace{-1.5cm}

\begin{center}
\small
\begin{tabular}{ll}
\textbf{Prepared for:} & Mr. Habib Fathallah \\
\textbf{Prepared by:} & The FSB Nexus Team \\
\textbf{Document type:} & Reference blueprint (supersedes prior planning documents) \\
\end{tabular}
\end{center}

\vspace{0.4cm}
\hrule
\vspace{0.4cm}

% ===========================================================================
\section{Executive Summary}

FSB Nexus is an intelligent assistant that helps students and prospective students
of the Faculty of Sciences of Bizerte (FSB) obtain reliable, sourced answers about
the faculty: its programs and licences, admission and administrative procedures,
course content, and general campus life. The system is built on a
Retrieval-Augmented Generation (RAG) architecture, ensuring that answers are grounded
in official documents and accompanied by citations rather than generated freely.

The foundation of the project is complete: authentication, the conversational
pipeline, the knowledge base, and the user interface are all implemented and
integrated end to end. This document summarizes what has been accomplished, presents
an assessment of the current state, and defines precisely what remains to be done and
how responsibilities are distributed across the team.

% ===========================================================================
\section{System Architecture}

The platform is composed of four cooperating layers:

\begin{itemize}
  \item \textbf{Frontend} --- a web application (Next.js) providing authentication,
        onboarding, the chat assistant, the profile, and the dashboard.
  \item \textbf{Backend} --- a REST API (FastAPI) handling authentication, user
        profiles, conversations, documents, and the bridge to the AI pipeline.
  \item \textbf{AI / RAG pipeline} --- a set of specialized agents that retrieve
        relevant document passages, generate a grounded answer, verify that the
        answer is supported by the sources, and attach citations.
  \item \textbf{Knowledge base} --- official FSB documents that are cleaned, split
        into passages, embedded, and stored in a vector database for semantic search.
\end{itemize}

Information flows as follows: the frontend calls the backend, the backend forwards a
question to the AI pipeline, the pipeline retrieves passages from the vector database,
produces a grounded answer with a language model, and returns the answer with its
sources to the user.

% ===========================================================================
\section{Work Accomplished}

The first phase of development is complete and merged into the main branch. The
following capabilities are implemented and functional.

\subsection{Backend}
\begin{itemize}
  \item Authentication with JSON Web Tokens: registration, login, and email
        verification, with hashed passwords and route protection.
  \item Profile management, including onboarding status and academic-year fields.
  \item Conversation endpoints that create sessions, store messages, and route
        questions to the AI pipeline, with per-user access control.
  \item Document endpoints for listing and administrative upload.
  \item A relational schema covering students, programs, departments, courses,
        documents, document passages, conversations, and messages.
\end{itemize}

\subsection{AI and Retrieval Pipeline}
\begin{itemize}
  \item Four specialized agents (Orientation, Academic, Administrative, and
        Learning), each with strict rules to answer only from provided sources and
        to cite the document and page.
  \item A complete pipeline: retrieval, relevance filtering, answer generation,
        citation formatting, and a groundedness check that blocks unsupported
        answers.
  \item A provider-agnostic language-model client, allowing the underlying model to
        be changed through configuration alone.
  \item A suite of automated tests covering prompts, citations, and the grading
        stages.
\end{itemize}

\subsection{Knowledge Base}
\begin{itemize}
  \item A corpus of official FSB documents, cleaned and organized by category.
  \item A structure-aware pipeline that splits documents into coherent passages
        while preserving titles, pages, and sources.
  \item A large set of passages embedded with a multilingual model and stored in the
        vector database, queryable by semantic similarity.
\end{itemize}

\subsection{Frontend}
\begin{itemize}
  \item Authentication pages, a dashboard, a chat interface, and a profile page.
  \item A responsive design with light and dark themes and a citation display.
\end{itemize}

\subsection{Integration}
\begin{itemize}
  \item The full stack has been assembled and verified locally: a user can register,
        log in, and receive sourced answers from the assistant.
\end{itemize}

% ===========================================================================
\section{Assessment of the Current State}

A careful review of the running system identified the areas that must be addressed to
move from a working foundation to a complete, dependable product.

\begin{itemize}
  \item \textbf{Conversational coverage.} The assistant is currently strict: it
        answers only when a directly matching passage is found and otherwise declines.
        It therefore does not yet handle greetings, general conversation, or
        broadly phrased questions, and it does not retain the context of a
        conversation across turns.
  \item \textbf{Retrieval quality.} Some answers are present in the knowledge base but
        are embedded within noisy source text, which prevents them from being
        retrieved reliably. Retrieval and source cleanliness both need to be
        strengthened.
  \item \textbf{Courses.} The assistant does not yet cover a student's own courses.
        Adding course content, and answering questions about it, is a core objective
        of the next phase.
  \item \textbf{Security and robustness.} Several hardening measures remain, including
        completing route protection on the client, rate limiting, password policy,
        stricter cross-origin configuration, and consistent error handling.
  \item \textbf{Experience and delivery.} Real-time (streaming) responses,
        conversation history, a fully data-driven dashboard, automated testing in
        continuous integration, and a reproducible deployment remain to be completed.
\end{itemize}

% ===========================================================================
\section{What Is Next: Objectives}

The next phase delivers a complete and professional product with the following
capabilities:

\begin{itemize}
  \item A chat assistant that answers any reasonable question about the faculty and
        its documents, understands general conversation, and stops declining valid
        questions, while never fabricating information.
  \item Course support, including the ability to bring in a student's own courses
        (with Google Classroom as the target integration) and to answer questions
        about their content.
  \item A polished, professional, and cohesive interface, with every screen connected
        to real data.
  \item A secure and robust platform suitable for a real deployment.
\end{itemize}

% ===========================================================================
\section{Roles and Responsibilities}

Work is organized into four roles, each owning a complete, demonstrable outcome. Roles
are described by responsibility and by the result they must deliver, not by a fixed
daily schedule. Assignment of team members to roles is decided by the team; pairing is
expected where a role naturally spans two areas.

\subsection{Role 1 --- Conversational Intelligence and Retrieval}
\textbf{Objective:} the assistant answers reliably and naturally, covering faculty
facts, document questions, and ordinary conversation, without fabricating information.

\textbf{Responsibilities:}
\begin{itemize}
  \item Introduce an intent-handling layer that distinguishes greetings, general
        conversation, and meta-questions (such as ``what can you do'') from factual
        questions, so that only factual questions follow the strict sourced path
        while the rest receive a helpful, bounded reply.
  \item Add short-term conversation memory so that follow-up questions are understood
        in context.
  \item Strengthen retrieval by combining keyword and semantic search and adding a
        re-ranking stage, so that answers present in the knowledge base are found.
  \item Refine the answering rules so the assistant provides what the sources support
        and clearly states what they do not, rather than declining outright.
  \item Improve the cleanliness of the most consulted source documents so that key
        information (programs, licences, procedures, location, contacts) is retrievable.
  \item Build an evaluation harness of representative questions that measures answer
        coverage and groundedness, serving as the shared quality benchmark for the
        whole team.
\end{itemize}

\textbf{Deliverable:} a measurable, significant increase in answered questions with
citations, while unsupported answers remain blocked.

\subsection{Role 2 --- Courses and Integrations}
\textbf{Objective:} a student can bring in their courses and ask questions about their
own course content.

\textbf{Responsibilities:}
\begin{itemize}
  \item Provide a reliable path for adding courses and uploading course materials, and
        store them in the appropriate tables.
  \item Ingest each student's course materials into the retrieval layer, scoped to
        that student, so that questions about a specific course are answered from the
        correct material.
  \item Expose endpoints to list courses, view a course and its materials, and refresh
        course content.
  \item Provide, together with Role 3, a courses screen that lists courses, shows their
        materials, and links directly into the assistant to ask about a course.
  \item Implement the Google Classroom integration (authorization, retrieval of the
        student's courses, coursework, and materials) as the target automated source.
        Because this integration depends on external authorization from Google and
        from the institution, a manual course-entry path is provided first so the
        capability is available regardless of that approval.
\end{itemize}

\textbf{Deliverable:} courses visible in the application and answerable in the
assistant, with Google Classroom as the automated source once authorized.

\subsection{Role 3 --- Frontend and Experience}
\textbf{Objective:} a professional, impressive, and cohesive interface in which every
feature is connected to real data.

\textbf{Responsibilities:}
\begin{itemize}
  \item Complete client-side route protection and session handling, with a proper
        redirect for unauthenticated or expired sessions and a working profile and
        logout.
  \item Deliver a real-time chat experience that displays answers as they are produced,
        with clear loading, error, and retry states, and safe rendering of formatted
        text.
  \item Present citations as an accessible source view that shows the document and page
        behind each answer.
  \item Provide a conversation-history panel to browse and resume previous
        conversations.
  \item Build the onboarding flow (student status, academic year, background, and
        interests) that personalizes the assistant.
  \item Replace all placeholder content on the dashboard with real data, including a
        summary of the student's courses and recent activity.
  \item Establish a consistent design system: unified spacing, typography, motion,
        responsiveness, light and dark themes, and accessibility.
\end{itemize}

\textbf{Deliverable:} a coherent, polished interface, free of placeholder content,
that presents the assistant and its features professionally.

\subsection{Role 4 --- Backend, Security, and Platform}
\textbf{Objective:} a secure, robust API that provides everything the other roles
require, and a dependable path to deployment.

\textbf{Responsibilities:}
\begin{itemize}
  \item Address the identified security items: request rate limiting, a password
        policy, stricter cross-origin configuration, consistent and transparent error
        handling, and safe handling of authentication tokens.
  \item Provide a streaming response endpoint so that the interface can display answers
        as they are generated.
  \item Provide conversation-history and feedback endpoints, persisting messages with
        their citations and capturing user feedback that feeds the evaluation benchmark.
  \item Make administrative document upload trigger ingestion into the retrieval layer.
  \item Support Role 2 with the server-side parts of the Google Classroom
        authorization and secure token storage.
  \item Establish reproducible operations: pinned dependencies, automated tests running
        in continuous integration, database migrations, and a straightforward
        deployment.
\end{itemize}

\textbf{Deliverable:} a hardened API exposing streaming, history, feedback, and
ingestion, with the platform prepared for deployment.

% ===========================================================================
\section{Shared Interfaces}

To allow the four roles to progress in parallel, the following interfaces are agreed
at the outset and each role builds against them:

\begin{itemize}
  \item \textbf{Streaming chat:} a conversation endpoint that accepts a message and an
        optional session identifier and emits the answer progressively, concluding with
        the session identifier and the list of citations.
  \item \textbf{Feedback:} an endpoint that records a rating and an optional note for a
        given answer.
  \item \textbf{Courses:} endpoints to list courses, retrieve a course with its
        materials, and refresh course content.
  \item \textbf{Pipeline call:} an entry point that accepts the message, the session,
        the student profile, the recent conversation history, and an optional course
        scope, and returns the answer, its citations, and the detected intent.
\end{itemize}

A feedback loop connects these interfaces: user feedback in the interface is recorded
by the backend, feeds the evaluation benchmark, highlights which sources to improve,
and thereby drives continuous improvement of answer quality.

% ===========================================================================
\section{Definition of Done}

The next phase is considered complete when:

\begin{itemize}
  \item Authentication is fully protected on both client and server, with rate limiting
        and a password policy in place.
  \item The assistant streams its answers, remembers the conversation, answers faculty
        and document questions with correct citations, and handles general
        conversation gracefully, with a measurable increase in answered questions and
        no unsupported claims.
  \item A student can add courses and receive answers about their own course content
        within the assistant.
  \item The interface is professional, cohesive, responsive, and free of placeholder
        content, with every screen connected to real data.
\end{itemize}

Additional production-grade objectives --- the fully authorized Google Classroom
integration, continuous integration and automated deployment, database migrations,
token refresh, password reset, and full internationalization and accessibility --- are
tracked as the subsequent set of improvements and are pursued as the core objectives
above are secured.

% ===========================================================================
\section{Conclusion}

The foundation of FSB Nexus is built and integrated. The next phase turns this
foundation into a complete, dependable, and professional assistant: one that answers
the questions students actually ask, incorporates their own courses, and presents a
polished experience on a secure platform. The roles and interfaces defined in this
document allow the team to advance in parallel toward that outcome, with a shared
evaluation benchmark keeping quality measurable at every step.

\end{document}
